\chapter*{Abstract}
\addcontentsline{toc}{chapter}{Abstract}
\begin{singlespace}
Childhood stunting has declined in Ghana, but the remaining burden is unevenly distributed. This study used the 2022 Ghana Demographic and Health Survey to examine how much residual variation in stunting lies between households and how much lies between survey communities. The analytic sample contained 4,928 children aged 0--59 months in 3,545 households and 616 communities. Survey-weighted stunting prevalence was 17.4\% (95\% confidence interval: 15.8\%--18.9\%). Bayesian hierarchical logistic models included household and community random intercepts and were extended sequentially with water and sanitation characteristics and maternal education. The final model was evaluated with eight chains and was challenged through prior, survey-weight, missing-data, age-functional-form, detailed-WASH, posterior-predictive, PSIS-LOO, and GEE sensitivity analyses. In the strengthened final model, the household random-effect standard deviation was 1.23 (95\% posterior interval: 0.94--1.53), compared with 0.39 (0.07--0.61) at the community level. The corresponding household variance partition coefficient was 0.31 and the community intraclass correlation was 0.03. Boys had higher posterior odds of stunting than girls (OR 1.50, 1.24--1.82), and children aged 24--35 months had higher odds than those aged 0--5 months (OR 2.91, 2.06--4.21). Unimproved drinking-water source was associated with higher odds of stunting (OR 1.50, 1.15--1.99), whereas the sanitation estimate was less clearly separated from one (OR 1.12, 0.86--1.46). Higher maternal education was associated with lower odds relative to no education (OR 0.37, 0.20--0.65). The main household-versus-community pattern remained across sensitivity analyses. The results are observational and point to substantial shared household conditions that are not captured by standard DHS covariates.
\end{singlespace}
\clearpage
